Service-level protection-mode mix of the 60 services of the reference
scenario before and after the selective QKD overlay. One row per mode that is
assigned to at least one service; in each row the upper, hatched light bar is
Phase~2 (PQC migration baseline, before any QKD overlay) and the lower, solid
dark bar is Phase~3 (after the selective QKD overlay has been planned and its
routes qualified, Figure~\ref{fig:qkd}). The number at the end of a bar is
the number of services in that mode, and a ``0'' next to the axis marks a
mode that no service uses in that phase. Modes are assigned per service, not
per physical link, by the policy engine at planning time (a qualified route
with an unconstrained key reserve is assumed; run-time reserve-based demotion
is not shown here). In the baseline, 50 services run PQ/T hybrid, 8 PQC only
and 2 traditional only, and no service uses QKD. The overlay changes the mode
of 12 services (20\% of the inventory), all of them high-assurance services
on the six qualified overlay links (core-0, core-3, metro0-2, metro2-2, dci-0
and dci-1): 8 IKEv2/IPsec and MACsec services move from PQ/T hybrid to
PQ/T + QKD, and the 4 OTN services on these links, whose encryptor profile
has no PQ/T hybrid capability, move from PQC only to QKD-assisted; the 4
services that remain in PQC only are the OTN services of the four core routes
that failed qualification (core-1, core-2, core-4 and core-5), and 42
services remain in PQ/T hybrid. The two standard-assurance services on
overlay links keep PQ/T hybrid because the QKD modes score lower for
standard-assurance services, and the 2 traditional-only services (the legacy
access endpoints acc-2 and acc-9, which have no PQC capability) are
unchanged. QKD therefore appears only after route qualification and only for
high-assurance services. Planning-model outputs of the seeded synthetic
scenario (topology seed 7, simulation seed 11); no carrier data are used.
